% Source: https://artofproblemsolving.com/community/c3h2200056_4_nt_problems_help
% Stable identity: aops-topic-2200056-q03
Suppose $n$ has $k+1$ digits, so $n=6\cdot10^k+m$, where $m$ (with $0\le m<10^k$) is the number left after removing the $6$. The condition $m=\dfrac {n}{25}$ gives
\[25m=6\cdot10^k+m,\qquad 24m=6\cdot10^k,\qquad m=\dfrac{10^k}{4}.\]
This is an integer only for $k\ge2$, and then $m=25\cdot10^{k-2}$, which indeed has $k$ digits and is less than $10^k$. Hence
\[n=6\cdot10^k+25\cdot10^{k-2}=625\cdot10^{k-2},\qquad k\ge2 .\]
The numbers are $\boxed{625,\ 6250,\ 62500,\ \dots}$, that is, $625\cdot10^j$ for $j=0,1,2,\dots$ (for example $625=25\cdot25$).
